% META: {"id": "1NSI-TC-AM-EXTRAIT", "chapitre": "1NSI-TYPES-CONSTRUITS", "type_objet": "amenagee", "capacites": ["1NSI-TYPES-CONSTRUITS-C1", "1NSI-TYPES-CONSTRUITS-C2", "1NSI-TYPES-CONSTRUITS-C5"], "mode_creation": "adaptation", "sources_inspiration": ["corpus_nsi/03_progressions/supports/premiere/P04/P04_version_amenagee_types_construits.md"], "status": "needs_review"}
\section*{Version aménagée --- Extrait Types construits}
\textit{Consignes séquencées, mise en page aérée, énoncés allégés.}

\bigskip

\subsection*{Exercice 1 --- Reconnaître un tuple (C1)}

Observe le programme ci-dessous :

\begin{python}
station = ("Tunis", 36.8, 10.2)
\end{python}

\bigskip

\textbf{Étape 1.} Complète le tableau :

\bigskip

\begin{center}
\begin{tabular}{|l|c|c|}
\hline
\textbf{Expression} & \textbf{Type} & \textbf{Valeur} \\
\hline
\lstinline{station[0]} & \dotfill & \dotfill \\
\hline
\lstinline{station[1]} & \dotfill & \dotfill \\
\hline
\lstinline{len(station)} & \dotfill & \dotfill \\
\hline
\end{tabular}
\end{center}

\bigskip

\textbf{Étape 2.} Essaie dans la console : \lstinline{station[1] = 37.0}

\medskip
Que se passe-t-il ? \dotfill

\medskip
Pourquoi ? Un tuple est \dotfill (mutable / non mutable).

\bigskip
\bigskip

\subsection*{Exercice 2 --- Créer une liste (C2)}

\textbf{Étape 1.} Crée un tableau de 4 prénoms :

\begin{python}
prenoms = [______, ______, ______, ______]
\end{python}

\textbf{Étape 2.} Ajoute un cinquième prénom :

\begin{python}
prenoms.______("Sami")
\end{python}

\textbf{Étape 3.} Combien d'éléments contient \lstinline{prenoms} maintenant ?

\medskip
\lstinline{len(prenoms)} = \dotfill

\bigskip
\bigskip

\subsection*{Exercice 3 --- Copier sans erreur (C5)}

\textbf{Lis} le programme ci-dessous, puis \textbf{entoure} la bonne réponse.

\begin{python}
notes = [12, 15, 9]
copie = notes
copie.append(18)
\end{python}

\bigskip

Que vaut \lstinline{notes} après ces trois lignes ?

\bigskip

\begin{center}
\fbox{\lstinline{[12, 15, 9]}} \qquad ou \qquad
\fbox{\lstinline{[12, 15, 9, 18]}}
\end{center}

\bigskip

Explique en une phrase : \dotfill

\medskip
\dotfill

\bigskip

\textbf{Corrige} la ligne 2 pour que \lstinline{notes} ne soit pas modifiée :

\begin{python}
copie = ______(notes)
\end{python}
